\clearpage
\begin{center}
{\itshape \large I - 19}
\end{center}
\vspace{1em}

\noindent\textbf{\textit{l}-ADIC REPRESENTATIONS} \hfill \textbf{I-19}

\noindent\rule{\textwidth}{0.4pt}

\noindent
and let $\mu$ be a Radon measure on $X$ (i.~e.~a continuous linear form on $C(X)$, \textit{cf.}~Bourbaki, \textit{Int.}, chap.~III, \S1). The sequence $(x_{\mathbf n})$ is said to be \underline{$\mu$-equidistributed}, or \underline{$\mu$-uniformly distributed}, if $\mu_{\mathbf n} \longrightarrow \mu$ \underline{weakly} as ${\mathbf n} \longrightarrow \infty$, i.~e.~if $\mu_{\mathbf n}(f) \longrightarrow \mu(f)$ as ${\mathbf n} \longrightarrow \infty$ for any $f \in C(X)$. Note that this implies that $\mu$ is positive and of total mass 1. Note also that $\mu_{\mathbf n}(f) \longrightarrow \mu(f)$ means that
\[
\mu(f) = \lim_{{\mathbf n} \to \infty} \frac{1}{n} \sum_{i=1}^{n} f(x_{\mathbf i}) \; .
\]

\noindent
LEMMA 1 - \underline{Let $(\phi_{\alpha})$ be a family of continuous functions on $X$ with the property that their linear combinations are dense in $C(X)$. Suppose that, for all $\alpha$, the sequence $(\mu_{\mathbf n}(\phi_{\alpha}))_{{\mathbf n}>1}$ has a limit. Then the sequence $(x_{\mathbf n})$ is equidistributed with respect to some measure $\mu$; \underline{it is the unique measure such that} $\mu(\phi_{\alpha}) = \lim\limits_{{\mathbf n} \to \infty} \mu_{\mathbf n}(\phi_{\alpha})$ \underline{for all} $\alpha$.}

If $f \in C(X)$, an argument using equicontinuity shows that the sequence $(\mu_{\mathbf n}(f))$ has a limit $\mu(f)$, which is continuous and linear in $f$; hence the lemma.

\noindent
PROPOSITION 1 - \underline{Suppose that $(x_{\mathbf n})$ is $\mu$-equidistributed. Let $\mathbf U$ be a subset of $X$ whose boundary has $\mu$-measure zero, and, for all $\mathbf n$, let $n_{\mathbf U}$ be the number of $\mathbf m \leq n$ such that $x_{\mathbf m} \in \mathbf U$. Then}

\hspace*{1.5em}$\lim\limits_{{\mathbf n} \to \infty} (n_{\mathbf U}/n) = \mu(\mathbf U)$.

Let $\mathbf U^0$ be the interior of $\mathbf U$. We have $\mu(\mathbf U^0) = \mu(\mathbf U)$. Let $\epsilon > 0$. By the definition of $\mu(\mathbf U^0)$ there is a continuous function $\phi \in C(X)$, $0 \leq \phi \leq 1$, with $\phi = 0$ on $X - \mathbf U^0$ and $\mu(\phi) \geq \mu(\mathbf U) - \epsilon$. Since $\mu_{\mathbf n}(\phi) \leq n_{\mathbf U}/n$ we have
